\documentclass[11pt]{article}

\usepackage[margin=1in]{geometry}
\usepackage{amsmath,amssymb,amsthm}
\usepackage{bm}
\usepackage{booktabs}
\usepackage{enumitem}
\usepackage[hidelinks]{hyperref}

\newtheorem{remark}{Remark}
\newtheorem{proposition}{Proposition}

\DeclareMathOperator*{\argmax}{arg\,max}
\DeclareMathOperator{\E}{\mathbb{E}}
\DeclareMathOperator{\Var}{Var}
\DeclareMathOperator{\Cov}{Cov}
\DeclareMathOperator{\tr}{tr}
\newcommand{\R}{\mathbb{R}}
\newcommand{\Norm}{\mathcal{N}}
\newcommand{\T}{^{\mathsf{T}}}
\newcommand{\given}{\,|\,}
\newcommand{\dd}{\mathrm{d}}
\newcommand{\half}{\tfrac{1}{2}}

\title{Innovation-Based Sequential Detection of Radio Transients:\\
Mathematical Foundations}
\author{KREMETART design notes}
\date{\today}

\begin{document}
\maketitle

\begin{abstract}
This note develops the statistical machinery underpinning the proposed transient
detector, deliberately separated from the imaging and software-pipeline design.
The detector treats each sky-pixel light curve as a Gaussian state-space model whose
quiescent dynamics are an integrated Wiener process. A Kalman filter ``whitens''
the time series into a sequence of independent normalised innovations; a transient
appears as a change in the distribution of those innovations, which we detect with
sequential change-detection statistics (CUSUM and its generalised-likelihood-ratio
extension). We define every quantity and acronym from first principles, give the
discrete-time process matrices in closed form, derive the innovation likelihood and
the marginal evidence, and characterise the operating point of the detector
(detection delay, false-alarm rate, multiple testing). A multi-scale bank of models
and its combination via Bayesian model averaging close the development.
\end{abstract}

\section{Overview and the detection problem}

We observe, at each sky pixel, a scalar light curve $\{y_k\}_{k=1}^{K}$ sampled at
times $t_k$, where $y_k$ is the (calibrated, source-subtracted) pixel flux in one
frame. In the absence of a transient the light curve is a slowly varying
``quiescent'' signal plus measurement noise. A \emph{transient} is a temporary
additive flux excess $a_k$ that switches on at an unknown time $t_0$ and persists for
an unknown duration with unknown amplitude and shape.

The detector has three layers:
\begin{enumerate}[leftmargin=*]
  \item A \emph{generative quiescent model} (Section~\ref{sec:iwp}): the light curve
        is an integrated Wiener process observed in noise. This encodes ``what normal
        looks like'' and provides a principled, recursive predictor.
  \item A \emph{whitening transform} (Section~\ref{sec:kf}): the Kalman filter turns
        the correlated light curve into a sequence of independent, unit-variance
        \emph{innovations} under the quiescent model. Detection reduces to testing
        whether those innovations remain zero-mean.
  \item A \emph{sequential test} (Section~\ref{sec:seqdet}): CUSUM / GLR statistics
        accumulate the signed innovation evidence and raise an alarm when it crosses
        a calibrated threshold, while remaining primed for a change at any future
        time.
\end{enumerate}
A bank of quiescent models at different temporal stiffnesses
(Section~\ref{sec:bank}) gives sensitivity across transient time-scales, in the
manner of a multi-scale matched filter.

\section{Notation and acronyms}

Throughout, lower-case bold denotes column vectors, upper-case bold matrices, and
$\Norm(\bm{m},\bm{\Sigma})$ the Gaussian density with mean $\bm{m}$ and covariance
$\bm{\Sigma}$. Table~\ref{tab:acro} collects the acronyms.

\begin{table}[h]
\centering
\small
\begin{tabular}{ll}
\toprule
Acronym & Meaning \\
\midrule
SDE   & Stochastic differential equation \\
GP    & Gaussian process \\
IWP   & Integrated Wiener process (the quiescent prior) \\
KF    & Kalman filter \\
LLR   & Log-likelihood ratio \\
LR / NP & Likelihood ratio / Neyman--Pearson (lemma) \\
SPRT  & Sequential probability ratio test (Wald) \\
CUSUM & Cumulative sum (change-detection) test (Page) \\
GLR   & Generalised likelihood ratio (test) \\
NIS   & Normalised innovation squared \\
NEES  & Normalised estimation error squared \\
ARL   & Average run length \\
ROC   & Receiver operating characteristic \\
FDR   & False discovery rate \\
BMA   & Bayesian model averaging \\
RFI   & Radio-frequency interference \\
\bottomrule
\end{tabular}
\caption{Acronyms used in this note.}
\label{tab:acro}
\end{table}

\section{The quiescent model: an integrated Wiener process}
\label{sec:iwp}

\paragraph{Continuous-time prior.}
A \emph{stochastic differential equation} (SDE) describes a continuous-time random
process through an ordinary differential equation driven by white noise. We model the
quiescent flux $f(t)$ as the $q$-times \emph{integrated Wiener process} (IWP): the
$q$-th derivative of $f$ is a Wiener process (Brownian motion). Collecting $f$ and its
first $q$ derivatives into a state vector
$\bm{x}(t)=\big(f(t),\,f^{(1)}(t),\dots,f^{(q)}(t)\big)\T\in\R^{q+1}$, the IWP is the
linear SDE
\begin{equation}
  \dd \bm{x}(t) \;=\; \bm{F}\,\bm{x}(t)\,\dd t \;+\; \bm{L}\,\dd\beta(t),
  \qquad
  \bm{F}=\begin{pmatrix} 0 & 1 & & \\ & 0 & \ddots & \\ & & \ddots & 1 \\ & & & 0\end{pmatrix},
  \quad
  \bm{L}=\begin{pmatrix}0\\\vdots\\0\\1\end{pmatrix},
  \label{eq:sde}
\end{equation}
where $\beta(t)$ is a scalar Wiener process with diffusion (spectral density)
$\sigma^2$. The scalar $\sigma^2$ is the \emph{driving variance}: it sets how rapidly
the top derivative wanders, hence how ``stiff'' or ``flexible'' the light curve is.
The IWP is a \emph{Gauss--Markov process}---Gaussian, and Markov in the state
$\bm{x}$---which is exactly what makes recursive (Kalman) inference exact and
$O(1)$ per sample.

\begin{remark}[Relation to Mat\'ern Gaussian processes and ``length scale'']
The IWP is non-stationary: its variance grows without bound, so it has no stationary
length scale. It is the natural smoothness prior of probabilistic numerics. A
\emph{stationary} alternative with an explicit length scale $\ell$ is the
Mat\'ern Gaussian process (GP) of half-integer smoothness $\nu=q+\tfrac12$, which
admits an SDE representation of the same order $q{+}1$ (the companion matrix $\bm{F}$
gains a stable characteristic polynomial set by $\ell$). In the bank of
Section~\ref{sec:bank} the ``scales'' are realised by varying $\sigma^2$ (and
optionally the order $q$) for the IWP, or $\ell$ for the Mat\'ern variant. We use
``stiffness'' for the qualitative knob: smaller $\sigma^2$ (or larger $\ell$) $=$
stiffer $=$ slower to follow changes.
\end{remark}

\paragraph{Exact discretisation.}
For a sampling interval $\Delta_k=t_k-t_{k-1}$, the SDE \eqref{eq:sde} integrates
\emph{exactly} to a discrete linear-Gaussian transition
\begin{equation}
  \bm{x}_k = \bm{A}(\Delta_k)\,\bm{x}_{k-1} + \bm{w}_k,
  \qquad \bm{w}_k\sim\Norm(\bm 0,\bm{Q}(\Delta_k)),
  \label{eq:disc}
\end{equation}
with $\bm{A}(\Delta)=\exp(\bm{F}\Delta)$ and
$\bm{Q}(\Delta)=\int_0^{\Delta}\!\exp(\bm{F}s)\,\bm{L}\sigma^2\bm{L}\T\exp(\bm{F}s)\T\dd s$.
Both are polynomials in $\Delta$. Indexing rows/columns by derivative order
$i,j\in\{0,\dots,q\}$,
\begin{equation}
  A_{ij}(\Delta)=\frac{\Delta^{\,j-i}}{(j-i)!}\ \ (j\ge i),\quad 0\ \text{otherwise};
  \qquad
  Q_{ij}(\Delta)=\sigma^2\,\frac{\Delta^{\,2q-i-j+1}}{(q-i)!\,(q-j)!\,(2q-i-j+1)} .
  \label{eq:AQ}
\end{equation}
The two cases used in practice are the constant-velocity ($q{=}1$) and
constant-acceleration ($q{=}2$) models:
\begin{equation*}
  q{=}1:\
  \bm{A}=\begin{pmatrix}1&\Delta\\0&1\end{pmatrix},\
  \bm{Q}=\sigma^2\begin{pmatrix}\Delta^3/3&\Delta^2/2\\\Delta^2/2&\Delta\end{pmatrix};
  \quad
  q{=}2:\
  \bm{A}=\begin{pmatrix}1&\Delta&\Delta^2/2\\0&1&\Delta\\0&0&1\end{pmatrix},\
  \bm{Q}=\sigma^2\!\begin{pmatrix}\tfrac{\Delta^5}{20}&\tfrac{\Delta^4}{8}&\tfrac{\Delta^3}{6}\\[2pt]
  \tfrac{\Delta^4}{8}&\tfrac{\Delta^3}{3}&\tfrac{\Delta^2}{2}\\[2pt]
  \tfrac{\Delta^3}{6}&\tfrac{\Delta^2}{2}&\Delta\end{pmatrix}.
\end{equation*}
For uniform sampling, $\bm{A}$ and $\bm{Q}$ are computed once per scale.

\section{Observation model and the Kalman filter}
\label{sec:kf}

We observe the flux (the zeroth derivative) in additive Gaussian noise:
\begin{equation}
  y_k = \bm{H}\bm{x}_k + v_k, \qquad \bm{H}=(1,0,\dots,0),\qquad v_k\sim\Norm(0,R_k),
\end{equation}
where $R_k$ is the per-frame noise variance, taken from the imaging weights / thermal
sensitivity (it may vary across pixels and frames). Together
\eqref{eq:disc} and this equation form a \emph{linear-Gaussian state-space model}, for
which the \emph{Kalman filter} (KF) gives the exact filtering distribution
$p(\bm{x}_k\given y_{1:k})=\Norm(\bm{x}_{k|k},\bm{P}_{k|k})$ recursively.

\paragraph{Recursion.} With $\bm{x}_{k|k-1}=\bm{A}\bm{x}_{k-1|k-1}$ and
$\bm{P}_{k|k-1}=\bm{A}\bm{P}_{k-1|k-1}\bm{A}\T+\bm{Q}$ (prediction):
\begin{align}
  e_k &= y_k - \bm{H}\bm{x}_{k|k-1}
      &&\text{(\emph{innovation}: one-step prediction error)},\\
  S_k &= \bm{H}\bm{P}_{k|k-1}\bm{H}\T + R_k
      &&\text{(\emph{innovation variance})},\\
  \bm{K}_k &= \bm{P}_{k|k-1}\bm{H}\T S_k^{-1}
      &&\text{(\emph{Kalman gain})},\\
  \bm{x}_{k|k} &= \bm{x}_{k|k-1} + \bm{K}_k e_k, \quad
  \bm{P}_{k|k} = (\bm{I}-\bm{K}_k\bm{H})\bm{P}_{k|k-1}
      &&\text{(\emph{update})}.
\end{align}
For numerical robustness across many parallel filters use the Joseph form
$\bm{P}_{k|k}=(\bm{I}-\bm{K}_k\bm{H})\bm{P}_{k|k-1}(\bm{I}-\bm{K}_k\bm{H})\T+\bm{K}_k R_k\bm{K}_k\T$,
or a square-root filter.

\paragraph{The whitening property.}
The decisive fact for detection is that, \emph{under the model}, the innovation
sequence is zero-mean and temporally white (Kailath's innovations approach):
\begin{equation}
  \E[e_k]=0,\qquad \E[e_k e_j] = S_k\,\delta_{kj}.
  \label{eq:white}
\end{equation}
The KF therefore acts as a \emph{whitening filter}: it removes all the predictable
(quiescent) structure, leaving a white residual. Anything that is \emph{not} explained
by the quiescent model survives in the innovations. Detection becomes a test on the
innovation sequence rather than on the raw, correlated light curve.

\section{The Bayesian evidence (prediction-error decomposition)}
\label{sec:evidence}

Because the innovations are independent Gaussians, the marginal likelihood of the data
under the model---the \emph{Bayesian evidence}---factorises as the product of the
one-step predictive densities (the ``prediction-error decomposition''):
\begin{equation}
  \log p(y_{1:K}\given \theta)
  = \sum_{k=1}^{K}\log\Norm(y_k; \bm{H}\bm{x}_{k|k-1}, S_k)
  = -\half\sum_{k=1}^{K}\Big[\log(2\pi S_k) + \tfrac{e_k^2}{S_k}\Big],
  \label{eq:evidence}
\end{equation}
where $\theta=(q,\sigma^2,\dots)$ are the model hyperparameters. This is the evidence
``that falls out of the Kalman filter.'' It serves two roles: (i) hyperparameter
inference, by maximising or marginalising \eqref{eq:evidence} over $\theta$ (the
objective is differentiable, so gradient methods or a grid both work); and (ii)
weighting the model bank in Section~\ref{sec:bank}. Note that for the linear-Gaussian
model \eqref{eq:evidence} is \emph{exact in closed form}; no particle/importance
sampling is required to evaluate or marginalise it.

\section{Innovation consistency: NIS and NEES}
\label{sec:nis}

Define the \emph{normalised innovation}
\begin{equation}
  z_k \;=\; \frac{e_k}{\sqrt{S_k}} \;\sim\; \Norm(0,1)\quad\text{under the quiescent model},
\end{equation}
and the \emph{normalised innovation squared} (NIS)
\begin{equation}
  \epsilon_k \;=\; e_k\T S_k^{-1} e_k \;=\; z_k^2 \;\sim\; \chi^2_{n_y},
\end{equation}
where $\chi^2_{n_y}$ is the chi-squared distribution with $n_y$ degrees of freedom and
$n_y=\dim(y_k)=1$ here. (The state-space analogue using the true state error,
$(\bm{x}_k-\bm{x}_{k|k})\T\bm{P}_{k|k}^{-1}(\bm{x}_k-\bm{x}_{k|k})\sim\chi^2_{n_x}$, is
the \emph{normalised estimation error squared}, NEES; it requires ground truth and is
used in simulation to verify filter consistency.)

A time-averaged NIS over $W$ samples, $\bar\epsilon=\tfrac1W\sum\epsilon_k$, has mean
$n_y$ and lies within $[\,\chi^2_{n_y W}(\tfrac\alpha2),\,\chi^2_{n_y W}(1-\tfrac\alpha2)\,]/W$
with probability $1-\alpha$ under the model. This is a \emph{consistency check}:
before trusting any detection, confirm the quiescent innovations are white and
unit-variance (NIS in bounds; flat autocorrelation of $z_k$). NIS is, however, a
two-sided single-sample statistic---it discards the sign of the innovation and does
not accumulate evidence---so it is a poor detector of faint or slow transients. The
detector proper accumulates the \emph{signed} innovations, as developed next.

\section{The transient as a change in distribution}
\label{sec:change}

Write the pixel flux as quiescent plus transient, $f_k=b_k+a_k$, with $a_k=0$ for
$k<t_0$. The KF is built to track $b_k$. By linearity of the filter, an additive input
$a_k$ injects a \emph{deterministic} mean into the innovation sequence:
\begin{equation}
  \mu_k \;=\; \E[e_k\given a_{1:k}] \;=\; a_k - \bm{H}\bm{A}\,\bm{x}^{a}_{k-1|k-1},
  \label{eq:mu}
\end{equation}
where $\bm{x}^{a}$ is the state estimate the filter would form when driven by $a$
alone. In words: the innovation sees the part of the transient that the filter has
\emph{not yet absorbed} into its state. The hypotheses are therefore
\begin{equation}
  H_0:\ z_k\sim\Norm(0,1)\quad\text{for all }k;
  \qquad
  H_1:\ z_k\sim\Norm(\delta_k,1)\ \text{for }k\ge t_0,\quad \delta_k=\mu_k/\sqrt{S_k}.
\end{equation}

\begin{remark}[Why stiffness controls slow-transient detectability]
For a step transient $a_k=A\,\mathbf{1}[k\ge t_0]$, equation \eqref{eq:mu} shows that
$\mu_{t_0}=A$ and that $\mu_k$ then decays toward zero as the filter ``catches up'',
with a time constant set by the closed-loop gain---hence by the stiffness. A
\emph{flexible} model (large $\sigma^2$) absorbs the step within a few samples,
$\delta_k\to0$ quickly, and a slow/faint transient becomes invisible: it is
\emph{explained away}. A \emph{stiff} model (small $\sigma^2$) lags, keeping
$\delta_k$ elevated over many samples, so the excess accumulates in the test
statistic. Slow transients thus require both a stiff quiescent reference and a long
integration horizon---two facets of the same requirement, and the motivation for the
multi-scale bank. The cost of excessive stiffness is false alarms on benign slow
drifts (gains, ionosphere, intrinsic variability).
\end{remark}

\section{Sequential detection}
\label{sec:seqdet}

\paragraph{Likelihood ratio and Neyman--Pearson (background).}
For deciding between two simple hypotheses, the Neyman--Pearson (NP) lemma states that
the most powerful test at a given false-alarm probability rejects $H_0$ when the
\emph{likelihood ratio} (LR) exceeds a threshold. For Gaussian innovations, the
per-sample \emph{log-likelihood ratio} (LLR) for a shift of size $\delta$ is
\begin{equation}
  \ell_k(\delta) = \log\frac{\Norm(z_k;\delta,1)}{\Norm(z_k;0,1)}
                 = \delta z_k - \half\delta^2 .
  \label{eq:llr}
\end{equation}

\paragraph{SPRT (Wald).}
The \emph{sequential probability ratio test} accumulates $\Lambda_k=\sum_{i\le k}\ell_i$
and stops as soon as $\Lambda_k\ge \log\frac{1-\beta}{\alpha}$ (declare $H_1$) or
$\Lambda_k\le\log\frac{\beta}{1-\alpha}$ (declare $H_0$), where $\alpha$ is the
type-I error probability (false alarm) and $\beta$ the type-II error probability
(miss). Among tests with these error rates, the SPRT minimises the expected number of
samples. Its limitation for our use is that it assumes a \emph{known} change time and
a single resolved decision; transient monitoring needs continuous surveillance with
an \emph{unknown} onset.

\paragraph{CUSUM (Page's test).}
The \emph{cumulative sum} test handles an unknown onset by resetting the accumulation
whenever the evidence turns non-positive:
\begin{equation}
  g_k = \max\!\big(0,\; g_{k-1}+\ell_k(\delta)\big),\qquad g_0=0,
  \qquad\text{alarm when } g_k\ge h .
  \label{eq:cusum}
\end{equation}
Substituting \eqref{eq:llr} gives the equivalent one-sided ``tabular'' form for a
positive (brightening) shift,
\begin{equation}
  g_k=\max\big(0,\;g_{k-1}+z_k-\kappa\big),\qquad \kappa=\delta/2,
\end{equation}
where the \emph{reference value} $\kappa$ is half the smallest standardised shift one
wishes to detect promptly. The $\max(0,\cdot)$ keeps the statistic primed for a change
beginning at any future sample; CUSUM is equivalent to an SPRT restarted at every
candidate change time. It is minimax optimal in Lorden's sense: it minimises the
worst-case expected detection delay subject to a lower bound on the mean time between
false alarms. A one-sided statistic (positive shifts only) is appropriate for
brightening transients and is roughly twice as sensitive as the two-sided version.

\paragraph{GLR (generalised likelihood ratio).}
CUSUM fixes the shift size $\delta$. When the amplitude is unknown, the
\emph{generalised likelihood ratio} test maximises the LLR over both the unknown onset
$t_0$ and the unknown amplitude. For a Gaussian mean shift the inner maximisation is
the maximum-likelihood estimate of $\delta$ over a candidate segment, with closed form
\begin{equation}
  G_k \;=\; \max_{1\le L\le N}\ \frac{1}{2L}\Bigg(\sum_{i=k-L+1}^{k} z_i\Bigg)^{\!2},
  \qquad\text{alarm when } G_k\ge h,
  \label{eq:glr}
\end{equation}
where $L=k-t_0+1$ is the (unknown) elapsed duration and $N$ is the maximum duration
searched. A short $L$ with large amplitude (a fast bright transient) and a long $L$
with small amplitude (a slow faint transient) are both admissible, so the GLR adapts
its integration length automatically. The maximising $L$ estimates the duration and
$\tfrac1L\sum z_i$ the standardised amplitude---useful science by-products. The
search window $N$ is precisely where a finite window legitimately enters the design:
\textbf{it is the longest transient the detector can optimally integrate}, not a
batching device. Lengthening $N$ improves slow-transient sensitivity at the cost of
computation and false-alarm exposure.

\paragraph{Matched-filter connection.}
If a transient temporal \emph{shape} $\bm{s}$ (e.g.\ an exponential flare or a
dispersed pulse) is known up to amplitude, the optimal statistic is the matched filter
$(\bm{s}\T\bm{z})^2/(\bm{s}\T\bm{s})$, which is $\chi^2_1$ under $H_0$. The GLR step
test \eqref{eq:glr} is the special case $\bm{s}=\text{boxcar}$; replacing the boxcar
with a template bank generalises the detector to specific light-curve morphologies.

\section{Operating characteristics and calibration}
\label{sec:arl}

\paragraph{Average run length.}
The operating point is summarised by two \emph{average run lengths} (ARLs):
$\mathrm{ARL}_0=\E[\text{samples to alarm}\given H_0]$ (mean time between false alarms;
to be made large) and $\mathrm{ARL}_1=\E[\text{detection delay}\given H_1]$ (to be made
small). The threshold $h$ trades them off; $\mathrm{ARL}_0$ grows essentially
exponentially in $h$. Closed-form approximations exist (Siegmund), but they assume
exactly white Gaussian innovations and are best used only to seed an empirical
calibration.

\paragraph{Receiver operating characteristic.}
The \emph{receiver operating characteristic} (ROC) plots detection probability against
false-alarm rate as $h$ varies. We obtain it by \emph{injection}: add synthetic
transients of known amplitude, duration and morphology into real
(source-subtracted) data and measure detections versus false alarms. This is also the
end-to-end test against an independent imager (e.g.\ DiSkO).

\paragraph{Multiple testing.}
With $P$ pixels each running a sequential test, the per-pixel threshold must be raised
to control the \emph{aggregate} false-alarm rate. For a target of one false alarm per
$T$ samples across the map, set the per-pixel $\mathrm{ARL}_0\approx P\,T$ (a
Bonferroni-type correction). When many candidate detections are expected, controlling
the \emph{false discovery rate} (FDR)---the expected fraction of declared detections
that are spurious---via the Benjamini--Hochberg procedure is less conservative than
controlling the family-wise error. Because the per-scale statistics
(Section~\ref{sec:bank}) are correlated, the multiplicity correction is best fixed by
empirical calibration on quiescent residuals rather than analytic union bounds.

\section{The multi-scale model bank and its combination}
\label{sec:bank}

Run $M$ quiescent models $\{\theta^{(m)}\}_{m=1}^{M}$ spanning a range of stiffnesses
(values of $\sigma^2$, optionally orders $q$). Each KF emits its own innovation
sequence $z_k^{(m)}$, evidence increment, and detection statistic $G_k^{(m)}$. Two
combinations are useful:

\paragraph{Bayesian model averaging (BMA) for the quiescent description.}
Treating the bank as a discrete hyperparameter prior, the posterior model weights
update recursively from the evidence increments \eqref{eq:evidence}:
\begin{equation}
  w_k^{(m)} \;\propto\; w_{k-1}^{(m)}\,
  \Norm\!\big(y_k;\,\bm{H}\bm{x}^{(m)}_{k|k-1},\,S_k^{(m)}\big),
  \qquad \sum_m w_k^{(m)}=1 .
\end{equation}
A forgetting factor (raising $w_{k-1}^{(m)}$ to a power $<1$ before re-normalising)
keeps the weights responsive to slow non-stationarity. \emph{Critically}, this
adaptation must run on a slower / more robust time-scale than detection: otherwise a
transient is absorbed by up-weighting a flexible model, hiding the very signal of
interest. In practice, freeze or heavily damp the weights during a candidate event.

\paragraph{Detection across scales.}
For declaring a transient, combine the per-scale detector outputs by the scale-wise
maximum, $G_k=\max_m G_k^{(m)}$, with the threshold raised to absorb the $M$-fold
multiplicity (again calibrated empirically). The argmax over scales reports the
best-matched transient time-scale. This realises a temporal multi-scale matched
filter: each scale is most sensitive to transients whose duration is comparable to,
or longer than, that model's tracking time.

\section{Practical considerations}

\begin{itemize}[leftmargin=*]
  \item \textbf{Initialisation.} Start each filter from a diffuse prior
  ($\bm{P}_{0|0}$ large); the first few $S_k$ are inflated, so suppress detections
  during a short warm-up.
  \item \textbf{Numerics.} Use Joseph-form or square-root covariance updates; with
  thousands of independent pixel$\times$scale filters of state dimension $q{+}1\le3$,
  the per-frame cost is dominated by small fixed-size linear algebra and vectorises
  cleanly.
  \item \textbf{Robustness to RFI.} Heavy-tailed outliers violate the Gaussian
  observation model and inflate false alarms. A Student-$t$ observation model or a
  simple innovation gate ($|z_k|$ clipped, with the sample down-weighted) restores
  robustness; the former makes the model non-Gaussian and is the one place a
  Rao--Blackwellised particle treatment (particles over an outlier/changepoint
  indicator, KF marginalising the continuous state) would be justified.
  \item \textbf{Validation order.} (i) Confirm filter consistency (NIS/whiteness) on
  quiescent data; (ii) calibrate $h$ for the target false-alarm rate on
  source-subtracted residuals; (iii) build ROC curves by injection; (iv) cross-check
  the recovered flux against an independent imager.
\end{itemize}

\section{Summary of the heuristic}

For each pixel and each scale $m$: run the IWP Kalman filter \eqref{eq:disc}, form the
normalised innovation $z_k^{(m)}$, and accumulate the one-sided GLR statistic
\eqref{eq:glr} over a window $N$ matched to the slowest transient of interest. Declare
a transient when $\max_m G_k^{(m)}$ exceeds a threshold calibrated for the desired
map-wide false-alarm (or false-discovery) rate, reporting the onset, duration and
amplitude from the maximising segment and the time-scale from the maximising model.
Quiescent-model adaptation (BMA weights, $\sigma^2$) runs on a slow, damped time-scale
so that transients are detected rather than absorbed.

\begin{thebibliography}{99}
\small
\bibitem{wald} A.~Wald, \emph{Sequential Analysis}, Wiley, 1947.
\bibitem{page} E.~S.~Page, ``Continuous inspection schemes,'' \emph{Biometrika}, 41(1/2):100--115, 1954.
\bibitem{lorden} G.~Lorden, ``Procedures for reacting to a change in distribution,'' \emph{Ann.\ Math.\ Statist.}, 42(6):1897--1908, 1971.
\bibitem{willsky} A.~S.~Willsky and H.~L.~Jones, ``A generalized likelihood ratio approach to the detection and estimation of jumps in linear systems,'' \emph{IEEE Trans.\ Autom.\ Control}, 21(1):108--112, 1976.
\bibitem{basseville} M.~Basseville and I.~V.~Nikiforov, \emph{Detection of Abrupt Changes: Theory and Application}, Prentice Hall, 1993.
\bibitem{siegmund} D.~Siegmund, \emph{Sequential Analysis: Tests and Confidence Intervals}, Springer, 1985.
\bibitem{kailath} T.~Kailath, ``An innovations approach to least-squares estimation, Part I,'' \emph{IEEE Trans.\ Autom.\ Control}, 13(6):646--655, 1968.
\bibitem{barshalom} Y.~Bar-Shalom, X.-R.~Li and T.~Kirubarajan, \emph{Estimation with Applications to Tracking and Navigation}, Wiley, 2001.
\bibitem{hartikainen} J.~Hartikainen and S.~S\"arkk\"a, ``Kalman filtering and smoothing solutions to temporal Gaussian process regression models,'' \emph{IEEE MLSP}, 2010.
\bibitem{sarkka} S.~S\"arkk\"a and A.~Solin, \emph{Applied Stochastic Differential Equations}, Cambridge University Press, 2019.
\bibitem{bh} Y.~Benjamini and Y.~Hochberg, ``Controlling the false discovery rate,'' \emph{J.\ R.\ Stat.\ Soc.\ B}, 57(1):289--300, 1995.
\end{thebibliography}

\end{document}
